\paragraph{Gruppo} Un gruppo è una coppia $\mathcal{G}=(\mathcal{X},\star)$, dove $\mathcal{G}$ è un insieme e $\star$ un'operazione binaria definita su $\mathcal{X}$. L'operazione $\star$ è:
\begin{itemize}
    \item associativa: $(a\star)\star c=a\star (b \star c)$, per qualche $a,b,c \in \mathcal{X}$
    \item $id \in \mathcal{X}$ è l'elemento \textit{identità} tale che: $a \star id=id \star a=a$ per qualche $a \in \mathcal{X}$
    \item $\forall a \in \mathcal{X} \ \exists b \in \mathcal{X}$ detta \textit{inversa} di $a$, tale che $a \star b = b\ star a= id$.
\end{itemize} 

\subparagraph{Gruppo Additivo} L'operazione $\star$ è l'addizione; $id$ è $0$; l'\textit{inversa} di $a$ è $-a$.

\subparagraph{Gruppo Moltiplicativo}L'operazione $\star$ è il prodotto; $id$ è $1$; l'\textit{inversa} di $a$ è $a^{-1}$.

\paragraph{Gruppo abeliano} Un gruppo $G=(\mathcal{X},\star)$ è abeliano se $\star$ è commutativa: $a \star b=b \star a$. 

\paragraph{Gruppo finito} Un gruppo $G=(\mathcal{X},\star)$ è finito se $\mathcal{X}$ è un insieme finito. 

\paragraph{Ordine di un gruppo} L'ordine di un gruppo $\mathcal{G}=(\mathcal{X},\star)$, scritto $ord(G)$, è il numero di elementi, ovvero $ord(G)=|\mathcal{X}|$. 


\paragraph{Gruppo additivo $\mathbb{Z}_n$} Sia $n\geq 2$. Il gruppo $\mathcal{G}=(\mathbb{Z}_n,+)$ è un gruppo abeliano finito di ordine $n$, con l'addizione $+$ modulo $n$. L'identità \textit{id} è $0$, l'\textit{inversa} di $a$ è $(-a) \pmod n$. Se $a=1, \dots, n-1$, allora $-a \pmod n=n-a$.

\paragraph{Gruppo moltiplicativo $\mathbb{Z}_p$} Sia $p\geq 2$ un primo. Sia $\mathbb{Z}_p^*=\mathbb{Z}\setminus\{0\}$. Il gruppo $\mathcal{G}=(\mathbb{Z}_p^
*,\cdot)$ è un gruppo abeliano finito di ordine $p-1$, col prodotto $\cdot$ modulo $p$. L'identità \textit{id} è $1$, l'\textit{inversa} di $a$ è $a^{-1}$. \\
$\mathbb{Z}_n^*=\mathbb{Z}_n \setminus \{ d \in \mathbb{Z}_n : \gcd(d,n) > 1\}$. \\ L'ordine di $\mathbb{Z}_n^*$ è $\phi(n)$. 

\paragraph{Funzione di Eulero} Per $n \geq 2$, $\phi(n)$ è il numero di interi positivi, minori di $n$, coprimi con $n$.  

\subparagraph{Teorema di Eulero} $\phi(n)$ può essere calcolato così: la \textbf{fattorizzazione} di $n$ è:
\[
n=\prod_{i=1}^{l}p_i^{e_i}, 
\] 

dove $p_i$ sono i primi, divisori distinti di $n$, ed $e_i \geq 1 $ per $i=1, \dots, l$. Allora 
\[
\phi(n)=\prod_{i=1}^{l}p_i^{e_i}-p_i(e_i-1)=n\prod_{p_i|n}(1-1/p_i).
\]

(Proprietà: $\phi(ab)=\phi(a)\phi(b)$, con $a,b$ divisori di $n$). 

\subparagraph{Corollario Teorema di Eulero} 
\begin{itemize}
    \item Se $p$ è primo, allora $\phi(p)=p-1$
    \item Se $p$ è primo, allora $\phi(p^e)=p^e-p^{e-1}$
    \item Se $p_1, \dots, p_l$ sono primi distinti, allora $\phi(\prod_{i=1}^{l}p_i)=\prod_{i=1}^{l}(p_i-1)$.
\end{itemize}

\paragraph{Ordine di un elemento} Per un gruppo finito $\mathcal{G}=(\mathcal{X}, \star)$, l'\textit{ordine} di un elemento $a \in \mathcal{X}$, scritto $ord(a)$, è il più piccolo intero positivo $i$, tale che 
\[
\underbrace{a \star a \star \dots \star a}_{i} = id.
\]

Se il gruppo è moltiplicativo: $a^i$; 

Se il gruppo è additivo: $i \cdot a$. Infatti in $(\mathbb{Z}_n, +)$ l'elemento $1$ ha ordine $n$. 

\paragraph{Teorema di Lagrange (ordine di un elemento)} \label{TeoLagrangeGruppi} Sia $\mathcal{G}=(\mathcal{X}, \star)$ un gruppo moltiplicativo di ordine $n$, e sia $g \in G$ un suo elemento. Allora l'ordine di ogni elemento $g$ divide l'ordine del gruppo $G$, ovvero \textbf{ord}(a) $|n$.

\begin{proof}
\end{proof}
% non viene dimostrato nello Stinson, su wikipedia usa concetto di classi laterali :|


\paragraph{Teorema} In un gruppo finito moltiplicativo $\mathcal{G}=(\mathcal{X}, \cdot)$, $\forall a \in \mathcal{X}$, l'ordine di $b=a^i$ è:
\[ 
\text{ord(b)}=\frac{\text{ord(a)}}{\gcd(\text{ord(a)}, i)}.
\]

\subparagraph{Corollario} Nel gruppo $\mathcal{G}=(\mathcal{X}, +)$, l'elemento $b$ ha ordine $n/\gcd(n,b)$.


\paragraph{Gruppo ciclico} Un gruppo abeliano finito $G=(\mathcal{X},\star)$ è un gruppo ciclico se esiste un elemento $a$ tale che ord(a)=$|\mathcal{X}|$. $a$ è detto \textbf{generatore} del gruppo. 

\vspace{0.2cm}
Infatti in $(\mathbb{Z}_n, +)$, con $n \geq 2$, l'elemento $1$ è un generatore. Poiché un elemento $a$ è un generatore se e solo se $\gcd(a,n)=1$, il numero di tutti i generatori di $(\mathbb{Z}_n, +)$ è $\phi(n)$.

Invece, $(\mathbb{Z}_p^*, \cdot)$, con $p \geq 2$ un primo, è un gruppo ciclico di ordine $(p-1)$ e un generatore di questo gruppo è detto \textbf{elemento primitivo} modulo $p$. Il numero di elementi primitivi è $\phi(p-1)=\phi(\phi(p))$. 

\paragraph{Teorema} Se $p$ è primo, allora $\mathbb{Z}_p^*$ è un gruppo ciclico. \\ In altre parole, $\exists g\in \mathbb{Z}_p^*$ generatore tale che $\mathbb{Z}_p^*=\{1,g,g^2,g^3,\dots,g^{p-2}\}$. \\ Se la fattorizzazione di $p-1$ è conosciuta allora è un semplice determinare l'ordine di ogni singolo elemento in $\mathbb{Z}_p^*$.

\paragraph{Sottogruppo} Sia $G=(\mathcal{X},\star)$ un gruppo finito e $\mathcal{Y} \in \mathcal{X}$. $H=(Y,\star)$ è un sottogruppo di G se anche H è un gruppo finito. 

\paragraph{Teorema} Sia G=(X,$\star$) un gruppo finito e sia Y $\in X$. Allora H=(Y,$\star$) è un sottogruppo di G $\iff$ è chiuso, ovvero $h_1 \star h_2 \in H \ \ \forall h_1, h_2 \in H$.  


\paragraph{Def} Sia $G=(\mathcal{X},\star)$ e sia $a \in \mathcal{X}$. Allora è facile vedere che il gruppo $(\langle a\rangle, \star)$ è un sottogruppo di G, dove l'ordine($\langle a\rangle$)$=$ord(a), ovvero l'ordine dell'elemento $a$ (per il gruppo X) coincide con l'ordine del (sottogruppo) gruppo generato da $a$.

\paragraph{Teorema di Lagrange (ordine di un sottogruppo)} Vista la definizione sopra, \ref{TeoLagrangeGruppi} si può riscrivere anche così:\\ Sia $H=(\mathcal{Y},\star)$ un sottogruppo di $G=(\mathcal{X},\star)$. Allora ord(H)$|$ord(G). 

\paragraph{Teorema Fondamentale sui gruppi ciclici} Sia $G=(\mathcal{X},\star)$ un gruppo finito di ordine $n$ e sia $a$ un generatore di G. Sia $d$ è un divisore di $n$. Allora, per ogni divisore $d$, esiste un unico sottogruppo di $G$ avente ordine $d$, ovvero:
\[
\qquad |\langle a^{n/d} \rangle| = d.
\]


%forse togliere tutti i \mathcal per gli insiemi
% fatta da pagina 528 a pagina 533 Stinson + aggiunti altri teo da fuori; fatta anche paragrafo 6.2.3 Stinson, riprendere per RSA direttamente (anche paragrafo prima )

